% Fact check of the CJAD 800 interview of August 20, 2025 with Stéphanie
% Valenzuela (then councillor for Darlington, candidate for borough mayor)
% on the Terrebonne bike path. Compile with pdflatex (twice), or run build.py.
\documentclass[11pt,letterpaper]{article}
\usepackage[T1]{fontenc}
\usepackage[utf8]{inputenc}
\usepackage[scaled=0.95]{helvet}
\renewcommand{\familydefault}{\sfdefault}
\usepackage[letterpaper,top=0.9in,bottom=0.95in,left=1in,right=1in]{geometry}
\usepackage{microtype}
\usepackage{enumitem}
\usepackage{titlesec}
\usepackage{fancyhdr}
\usepackage{xcolor}
\usepackage{booktabs}
\usepackage{longtable}
\usepackage{array}
\usepackage{textcomp}
\usepackage[hidelinks]{hyperref}
\urlstyle{same}
\newcommand{\docversion}{2026-10-26-12-38-56} % set by build.py (America/Toronto)
\hypersetup{pdfsubject={version \docversion}}

\setlength{\parindent}{0pt}
\setlength{\parskip}{0.5em}
\linespread{1.03}

\titleformat{\section}{\normalfont\large\bfseries}{\thesection}{0.6em}{}[\vspace{1pt}\titlerule]
\titlespacing*{\section}{0pt}{1.1em}{0.4em}
\titleformat{\subsection}{\normalfont\bfseries}{\thesubsection}{0.6em}{}
\titlespacing*{\subsection}{0pt}{0.8em}{0.25em}

\pagestyle{fancy}
\fancyhf{}
\renewcommand{\headrulewidth}{0pt}
\fancyfoot[C]{\footnotesize\color{gray}CJAD 800, August 20, 2025 -- fact check -- page \thepage}
\fancyfoot[R]{\tiny\color{gray}v.\,\docversion}

\newcommand{\repourl}{https://github.com/OpenMobilityData/Terrebonne-open-sources}
\newcommand{\repo}[1]{\href{\repourl}{#1}}
\newcommand{\doc}[2]{\href{\repourl/blob/main/#1}{#2}}
\newcommand{\arurl}{https://github.com/OpenMobilityData/Terrebonne-access-requests}
\newcommand{\pagesurl}{https://openmobilitydata.github.io/Terrebonne-access-requests}
\newcommand{\cliptime}[1]{\textsf{\small[#1]}}
\newcommand{\pgl}[2]{\href{\pagesurl/2026-581/raw/01_emails_mandate_counts.pdf\#page=#1}{\textsf{\small[2026-581, p.~#2]}}}
\newcommand{\pg}[1]{\pgl{#1}{#1}}
\newcommand{\exppg}[1]{\href{\repourl/blob/main/2023-exp-mobility-study/raw/EXP_Etude_de_mobilite_rue_de_Terrebonne_2023.pdf}{\textsf{\small[EXP 2023, p.~#1]}}}
\newcommand{\wc}[1]{\href{\repourl/blob/main/2020-2021-working-committee/raw/Comite_de_travail_recommandations_2021-06-13.pdf}{\textsf{\small[committee note, June 13, 2021, p.~#1]}}}
\newcommand{\quot}[1]{\begin{quote}\small #1\end{quote}}

\begin{document}

\begin{center}
{\Large\bfseries CJAD 800, August 20, 2025: the candidate's statements on the Terrebonne bike path checked against the record}\\[0.4em]
{\large Fact check of the interview with Stéphanie Valenzuela}\\[0.6em]
Richard Hoge, Côte-des-Neiges--Notre-Dame-de-Grâce -- October 26, 2026\\
Published with the transcript at \repo{\repourl}
\end{center}

\section*{Summary}

On August~20, 2025, CJAD~800's afternoon show interviewed Stéphanie
Valenzuela, then city councillor for the Darlington district and Ensemble
Montréal's candidate for borough mayor of Côte-des-Neiges--Notre-Dame-de-Grâce
(elected November~2, 2025), about her priorities, the Terrebonne bike path and
a McGill study published the week before. Each statement of fact about the
bike path was compared with the public record. In short:

\begin{enumerate}[nosep]
\item \textbf{The working committee's recommendations were not « completely disregarded ».} The committee's mandate (resolution CA20~170210, 2020) was to recommend improvements to mobility on rue de Terrebonne, not to design a bike path. Its unanimous recommendation (note of June~13, 2021) was that the City \emph{study} scenario~5\,bis, a one-way eastbound Terrebonne between Girouard and Grand, Madison or Benny, together with the Somerled scenarios; it adopted no layout and was itself divided on the western half. The City then commissioned the EXP mobility study (counts in October 2022, report published October~30, 2023), and the 2024 project made Terrebonne one-way. Where the project follows the committee (a one-way street, studied first) and where it goes beyond it (the whole length to Belmore; protected lanes rather than the committee's 4.1~m shared lane) is in the documents. The candidate was not a member of the committee, which the councillor for Notre-Dame-de-Grâce chaired.
\item \textbf{« Implemented … without consulting the residents » leaves out what the record shows.} Between the 2020 lanes and the 2024 project: the working committee's eight meetings with institutions, school boards and residents (November 2020 to May 2021); the EXP study and its public release of October~30, 2023; a public information session on November~30, 2023; the borough council's vote of February~5, 2024 (3--2, the candidate voting against) and the City council's of February~19, 2024. Whether that consultation was sufficient is a judgement; that there was none is not what the documents show.
\item \textbf{The 2020 history is stated correctly:} temporary lanes installed in summer 2020 under the previous administration and removed that September (motion of September~8, 2020).
\item \textbf{The audit.} In August 2025 the candidate described her party's proposed audit as looking at the path's « security, … implementation, … impact on local residents ». What followed is on the record: the City's own audit of the cycling network, launched February~10, 2026 under the principle « recueillir les données, analyser et corriger au besoin », with results due at the end of 2026 and Terrebonne named; and, in the borough, a process documented by access request 2026-581 in which the elected officials' objective of restoring the parking strip was set on June~10, 2026 before any study, the City's street-design service found « plusieurs compromis importants sur le plan de la sécurité » in the borough's options (June~22), no parking or resident data were collected, and the winter removal of one direction was announced on October~6, 2026, before the City's audit reported.
\item \textbf{The McGill study is described in terms it does not use.} The study (Romm \emph{et al.}, \emph{Journal of Cycling and Micromobility Research}, 2025, open access) measures road space per traveller by mode and borough and models what doubling the micromobility lanes would do to that allocation; it finds that even then street space « remains biased in favour of automobiles ». It does not claim to measure « the actual reality of Montrealers when it comes to cycling », and it states its own limits (2018 trip data, borough-scale averages, multi-use paths excluded). The candidate gives no basis for calling it not comprehensive; the host's paraphrase, « double … and it still wouldn't be enough », is close to the study's conclusion.
\item \textbf{« Put in different places where it's really questionable … a safe place », and cyclists who « constantly » report blind spots: no place and no source are given,} so the statements cannot be verified. For Terrebonne, the 2024 layout was designed through the EXP study; the City's safety concerns in the 2026 file concern the borough's 2026 options for restoring parking, not the 2024 layout; and the borough's counter recorded 135,276 bicycle passages at Kensington in the lane's first full year of counting (July~26, 2025 to July~25, 2026).
\item \textbf{Framing, noted without judgement:} the needs listed are « bike paths or … your vehicle or … public transportation », and the audit is justified because « there \emph{are} people that use … scooters or bicycles » (stress as heard); the host calls Côte-des-Neiges and NDG « boroughs » and is not corrected (they are the two halves of one borough).
\end{enumerate}

\section{What was checked, and how}

\textbf{The interview.} CJAD~800 (Bell Media), afternoon show, clip
« Ensemble Montréal's Stephanie Valenzuela is running for CDN-NDG borough
mayor », published August~20, 2025, 8~min~28~s; host Aaron Rand. The
transcript published with this document was drafted by machine and then
checked passage by passage against the audio.

\textbf{Sources.} The 2020 council motions and the working committee's note
of June~13, 2021 (\doc{2020-2021-working-committee/}{this repository}); the
2023 EXP mobility study; the 2024 council records (CA24~170018, CM24~0179,
CA24~170131); the City's release of February~10, 2026 on the cycling-network
audit; access request 2026-581 (borough response of October~9, 2026,
bundle~01, cited by page); the Kensington counter's daily table (2026-581,
pp.~74--81); the McGill study (Romm \emph{et al.}, 2025); the 2021 Census.

\textbf{Terms.} \emph{Consistent} means the record supports the statement as
made; \emph{cannot be verified} that the record does not bear on it. A
statement can be accurate in its terms and still leave out what the record
also shows; the summary says what was left out.

\section{Findings in detail}

\subsection{The Terrebonne history and the working committee}

\cliptime{03:34} « If anybody knows the history of Terrebonne, in the last
mandate, under the last administration, they had implemented it and then they
had removed it. »

\textbf{Consistent.} Temporary lanes were installed on Terrebonne in the
summer of 2020 under the « rues, secteurs et corridors à vocation familiale
et active » programme (council file of June~22, 2020) and removed after the
motion of September~8, 2020 (\doc{2020-temporary-lanes/}{documents}).

\cliptime{03:34} « And they created this working group where you had
institutions and people that were part of this discussion to see how can we
put a bike path and how can we get it to work for the residents and for the
people that want to use it. And it was completely, you know, disregarded when
they decided to re-implement a bike path. »

\textbf{Not what the record shows.} Resolution CA20~170210 (October~5, 2020)
created « un groupe de travail … ayant pour mandat de proposer des
recommandations qui visent l'amélioration de la mobilité sur la rue de
Terrebonne », chaired by the councillor for Notre-Dame-de-Grâce, with borough
staff, the school boards, the street's religious institutions, the HLM and
residents; its mandate was mobility on the street, not « how can we put a bike
path ». The committee met eight times from November~19, 2020 to May~19, 2021.
Its note of June~13, 2021 \wc{5} recommends unanimously that the council ask
the City to study scenario~5\,bis (rue de Terrebonne one-way eastbound between
Girouard and, to the west, Grand, Madison or Benny, with the effect on
neighbouring streets taken into account), « en conjonction et en complément »
with the Somerled scenarios; by majority, that scenario~1 also be studied.
Three members had serious reservations about any change west of Grand; three
others and the chair wanted the lane to reach Benny. The committee recommended
a study; it adopted no layout.

What followed: the City commissioned EXP's mobility study, with counts at 17
intersections in October 2022 and a report published on October~30, 2023
(\doc{2023-exp-mobility-study/}{documents}); the borough council approved the
2024 project on February~5, 2024 (CA24~170018, 3--2) and the City council on
February~19, 2024 (CM24~0179); the ordinance of June~3, 2024 (OCA24~17042)
made Terrebonne one-way « en vue d'ajouter un aménagement cyclable protégé »
(\doc{2024-project-approval/}{documents}). The 2024 project follows the
committee in making the street one-way after a study, and goes beyond it in
extent (the whole length from Girouard to Belmore) and in form (protected
one-way lanes on both sides rather than the 4.1~m shared lane of scenario
5\,bis). The candidate was not a member of the committee.

\subsection{« Without consulting the residents »}

\cliptime{03:34} « And the real issue here is the way that the administration
went about and implemented it without consulting the residents. »

\textbf{Leaves out what the record shows.} The steps listed above: the
committee (institutions, school boards, residents), the study and its public
release, a public information session on November~30, 2023 (recorded in the
borough's decisional file 1246880002), and two council votes. The candidate
voted against the project at the borough council on February~5, 2024. Whether
those steps amounted to adequate consultation is a judgement on which the
record is silent.

\subsection{The audit}

\cliptime{03:34} « So there's a lot of work that needs to be done in the audit
in terms of looking at its security, looking at its implementation, looking at
its impact on local residents. And that's where we stand when it comes to at
least the bike path on Terrebonne. » \cliptime{04:46} « … that's why we go
back to the fact that we do need this audit … »

\textbf{A statement of intent at the time; what followed is on the record.}
The City's audit of the cycling network was launched by the new
administration on February~10, 2026: a technical analysis led by the Service
de l'urbanisme et de la mobilité, data collection and field observation
through the peak season, results expected at the end of 2026, Terrebonne named
among the streets where the City works with the borough on safety
(\doc{2026-city-audit-and-plan-velo/}{release and timeline}). In the borough,
the file released under access request 2026-581 shows: on June~10, 2026 the
borough's engineering division wrote to the City that the elected officials
wished to revisit the 2024 layout « avec l'objectif de redonner la bande de
stationnement qui avait été retranchée » \pg{8}; the City's street-design
service returned an express analysis on June~17, which the borough's own
summary of June~22 describes as finding that the options « présentent
plusieurs compromis importants sur le plan de la sécurité, du confort des
usagers vulnérables et de la cohérence du réseau cyclable » \pg{6}; the
analysis itself was withheld; no parking inventory, occupancy count or
resident survey appears in the 508 pages; and on October~6, 2026 the borough
announced the winter removal of one direction, before the City's audit
reported. Looking at « security, implementation, impact on local residents »
is what the audit was said to be for; the file shows the borough fixing the
parking objective first and collecting no data on those three points.

\subsection{The McGill study}

\cliptime{04:46} Host: « that study from McGill … suggested Montreal could, in
fact, double the number of bike paths we have, and it still wouldn't be
enough. » Candidate: « I find it very interesting. You know, there's a lot of
studies that come out that come in favor or that come against the
implementation of bike paths. … I understand that the study came out
recently, but I don't think it's a comprehensive study on the actual reality
of Montrealers when it comes to cycling. »

\textbf{The study, as published.} Daniel Romm, José Arturo Jasso Chávez, Lexi
Kinman, Pegah Salsabilian, Grant McKenzie and Kevin Manaugh (Department of
Geography, McGill University), « The cars are going to be alright: Examining
micromobility infrastructure space allocation and potential improvement
scenarios in Montréal », \emph{Journal of Cycling and Micromobility Research}
5 (2025) 100071, \url{https://doi.org/10.1016/j.jcmr.2025.100071}, online
June~7, 2025, open access (CC~BY-NC). It measures, per borough, the road space
allocated to automobiles and to micromobility per traveller of each mode (trip
data from the 2018 origin--destination survey, infrastructure of 2022), defines
an Equal Infrastructure Allocation score with parity at zero, and models two
scenarios: doubling the Bixi stations and doubling the area of dedicated
micromobility lanes. Its conclusion: « even with MI improvements at a scale
dramatically larger than anything now proposed, such as doubling all existing
dedicated micromobility lanes, street space allocation remains biased in
favour of automobiles » (p.~12), with « almost no discernable effect » on
automobile space (p.~9). The host's paraphrase is close to that sentence. The
study's own limitations section (4.4) lists borough-scale averages, 2018 trips
against 2022 infrastructure, no separation of bus space, commercial transport
omitted and multi-use paths excluded.

\textbf{The candidate's assessment cannot be verified.} The study does not
set out to measure « the actual reality of Montrealers when it comes to
cycling »; it measures space allocation. No basis is given for calling it not
comprehensive. « A lot of studies … in favor or … against the implementation
of bike paths » is said of a study that takes no position on whether bike
paths should exist. Press coverage: McGill Newsroom, « Study finds Montreal
cycling infrastructure doesn't match demand »; CTV News Montreal, August~18,
2025, « Should Montreal double its number of bike paths? One study says
`yes' » (micromobility users about 5\,\% of commuters, their infrastructure
about 2\,\% of road space).

\subsection{Safety of the bike paths}

\cliptime{04:46} « But the bike paths that have been put under this
administration have been put in different places where it's really
questionable to what extent it's actually a safe place to put them in. There's
cyclists that tell us constantly that we don't use that bike path because it's
actually really not safe. There's a lot of blind spots when you use that bike
path. There's actually a lot of drivers that don't see us when we're passing
by there. »

\textbf{Cannot be verified:} no place, count or source is named. For
Terrebonne, the 2024 layout followed the EXP study's design work
\exppg{51}; the City's expressed safety concerns in the 2026 file bear on the
borough's 2026 options for restoring parking (a designated roadway and a
narrow bidirectional lane), not on the 2024 layout \pg{6}; the borough's
counter at Kensington recorded 135,276 bicycle passages in the year from
July~26, 2025 (2026-581, pp.~74--81), and the borough's own Plan local de
déplacements 2024--2034, adopted unanimously on May~6, 2024 with the candidate
present and voting, treats protected lanes as a safety measure for vulnerable
users (\doc{2024-cdn-ndg-pld/}{documents}).

\subsection{Other statements}

\cliptime{01:44} « it's the most populous borough on the island of Montreal »:
\textbf{consistent}; Côte-des-Neiges--Notre-Dame-de-Grâce is the most
populous of the City's nineteen boroughs (2021 Census).

\cliptime{01:44} Taxes « increased … significantly », services « diminishing »,
rats, parks, permits, a tree « earmarked to be cut … four years »: outside the
scope of this check (not about the bike path); not examined.

\cliptime{06:10} Host: « in the Côte-des-Neiges borough … in the NDG
borough »: Côte-des-Neiges and Notre-Dame-de-Grâce are the two halves of one
borough, with five districts (Darlington, Côte-des-Neiges, Snowdon,
Notre-Dame-de-Grâce, Loyola); the candidate answers without correcting the
host. Noted without judgement.

\subsection{Framing}

Noted as observations on wording, not on intent. \cliptime{04:46} The
mobility needs listed are « using bike paths or, you know, your vehicle or
using public transportation »; the sentence addresses the listener as a car
owner and leaves out households without a car and people who walk. City-wide
in 2016, by main mode of commuting: 367,533 car drivers, 22,656 passengers,
274,575 public transit, 57,390 walking, 30,278 bicycle (Census). \cliptime{04:46}
« there \emph{are} people that use, you know, scooters or bicycles »: the
stress, as heard on the audio, falls on « are ».

\section{What followed}

The candidate was elected borough mayor on November~2, 2025. On October~6,
2026 the borough announced its orientation to convert one direction of the
Terrebonne lane to parking each winter, with the vote expected on December~7,
2026; the file behind it is the subject of access request 2026-581
(\href{\arurl/tree/main/2026-581}{summaries and documents}). Her later radio
statements (CBC \emph{Daybreak}, November~6, 2025 and May~20, 2026) are being
transcribed and will be checked separately.

\section{What this check does not establish}

It compares statements with records. It does not establish why any statement
was made, and it draws no conclusion about intent. The City's analysis of June
2026 and the consultant's note are withheld; the City's audit had not reported
when this check was written.

\section{Sources}

\begin{small}
\begin{longtable}{>{\raggedright\arraybackslash}p{4.2cm}>{\raggedright\arraybackslash}p{10.6cm}}
\toprule
Source & Where \\
\midrule
\endhead
Interview & CJAD~800, August~20, 2025: \href{https://omny.fm/shows/cjad-800/ensemble-montr-al-announces-their-c-te-des-neiges-notre-dame-de-gr-ce-team-headed-by-stephanie-valenzuela}{Omny clip page} (Montreal Now with Aaron Rand feed); transcript: \doc{2025-08-20-cjad-valenzuela/derived/2025-08-20_CJAD_Valenzuela_transcript_en.md}{this folder} \\
2020 lanes & \doc{2020-temporary-lanes/}{Council file of June~22, 2020; motion of September~8, 2020} \\
Working committee & \doc{2020-2021-working-committee/}{Resolution CA20~170210; note of June~13, 2021; file 1215284005; vote results} \\
EXP 2023 & \doc{2023-exp-mobility-study/raw/EXP_Etude_de_mobilite_rue_de_Terrebonne_2023.pdf}{Étude de mobilité -- rue de Terrebonne} (p.~51 on the bidirectional option) \\
2024 approvals & \doc{2024-project-approval/}{CA24~170018, CM24~0179, CA24~170131}; public information session of November~30, 2023 (file 1246880002) \\
PLD 2024--2034 & \doc{2024-cdn-ndg-pld/}{Plan local de déplacements} (CA24~170089, May~6, 2024) \\
City audit & \doc{2026-city-audit-and-plan-velo/}{Release of February~10, 2026 and the Réalisons Montréal timeline} \\
Request 2026-581 & Borough response of October~9, 2026, bundle 01: \href{\pagesurl/2026-581/raw/01_emails_mandate_counts.pdf}{PDF} (pp.~6, 8, 74--81); summaries in the \href{\arurl/tree/main/2026-581}{access-requests repository} \\
McGill study & Romm \emph{et al.}, 2025, \url{https://doi.org/10.1016/j.jcmr.2025.100071}; \href{https://www.mcgill.ca/newsroom/channels/news/study-finds-montreal-cycling-infrastructure-doesnt-match-demand-366423}{McGill Newsroom release}; \href{https://www.ctvnews.ca/montreal/article/montreal-should-double-its-number-of-bike-paths-study/}{CTV News Montreal, August~18, 2025} \\
Census & Statistics Canada, Census of Population 2016 and 2021, Ville de Montréal, main mode of commuting; borough populations \\
\bottomrule
\end{longtable}
\end{small}

\end{document}
